%
% device-table.tex
%
% Z specification of PortAudio's device table as voxd keeps it
% (src/punt_vox/voxd/conversation_mode/device_table.py): the count of
% streams reserved on the table, whether the library is initialized, and
% the three properties the tests in tests/conversation_mode/ assert by name.
%

\documentclass[a4paper,10pt,fleqn]{article}
\usepackage[margin=1in]{geometry}
\usepackage{fuzz}

\begin{document}

\title{The PortAudio Device Table: A Z Specification}
\author{Punt Labs}
\date{November 2026}
\maketitle

\section{Introduction}

PortAudio enumerates the system's audio devices once, when it is
initialized, and answers every later query from that copy. The only way to
read the table again is to terminate the library and initialize it again,
and terminating closes every stream still open. voxd therefore keeps one
\texttt{DeviceTable} per process: it counts the streams reserved on the
current table and re-reads the table only when that count is zero.

A stream is counted from before its native stream exists. The opener takes
a reservation under the table's lock, then creates the stream; the
reservation is released once, when the stream closes or when creating or
starting it fails. The model names three properties, and the tests carry
the same names:

\begin{enumerate}
  \item[P1] \emph{Refresh never terminates PortAudio while a stream is
    reserved}: a refresh with the count above zero keeps the table.
  \item[P2] \emph{Refresh terminates PortAudio only while it is
    initialized}: an initialize that fails leaves the library terminated,
    and the next refresh initializes directly instead of terminating again.
  \item[P3] \emph{Reservations balance}: every release answers one
    reservation; releasing more than were reserved is a bug.
\end{enumerate}

\section{Types}

\begin{zed}
  ZBOOL ::= ztrue | zfalse
\end{zed}

\section{State}

\begin{schema}{DeviceTable}
  openStreams : \nat \\
  initialized : ZBOOL
\end{schema}

\texttt{openStreams} is the number of reservations held, each for one
stream that exists or is about to. \texttt{initialized} is whether PortAudio
is initialized; it is \texttt{zfalse} only between a terminate and the
initialize that failed to follow it.

\begin{schema}{InitDeviceTable}
  DeviceTable'
\where
  openStreams' = 0 \\
  initialized' = ztrue
\end{schema}

Importing \texttt{sounddevice} initializes PortAudio, and a table is first
consulted after that import.

\section{Reservations}

\begin{schema}{Reserve}
  \Delta DeviceTable
\where
  openStreams' = openStreams + 1 \\
  initialized' = initialized
\end{schema}

\begin{schema}{Release}
  \Delta DeviceTable
\where
  openStreams \geq 1 \\
  openStreams' = openStreams - 1 \\
  initialized' = initialized
\end{schema}

The precondition of \texttt{Release} is P3: the implementation raises
\texttt{ValueError} when no stream is reserved. A \texttt{Reservation}
object applies \texttt{Release} at most once, so a stream closed twice
releases once.

\section{Refresh}

\texttt{terminated!} reports whether the refresh terminated the library
before initializing it.

\begin{schema}{Refresh}
  \Delta DeviceTable \\
  terminated! : ZBOOL
\where
  openStreams = 0 \\
  terminated! = initialized \\
  openStreams' = openStreams \\
  initialized' = ztrue
\end{schema}

\begin{schema}{RefreshKept}
  \Xi DeviceTable \\
  terminated! : ZBOOL
\where
  openStreams \geq 1 \\
  terminated! = zfalse
\end{schema}

\texttt{Refresh} and \texttt{RefreshKept} together are total over
\texttt{openStreams}, and only \texttt{Refresh} terminates: P1. In
\texttt{Refresh}, \texttt{terminated!} equals \texttt{initialized}: P2, the
library is terminated only while it is initialized.

\begin{schema}{FailedInitialize}
  \Delta DeviceTable \\
  terminated! : ZBOOL
\where
  openStreams = 0 \\
  terminated! = initialized \\
  openStreams' = openStreams \\
  initialized' = zfalse
\end{schema}

A \texttt{FailedInitialize} is a \texttt{Refresh} whose initialize step
raised. It leaves \texttt{initialized} false, so the \texttt{Refresh} that
follows has \texttt{terminated!} = \texttt{zfalse}: it initializes directly
and re-reads the table, which is what the regression test for P2 asserts.
Without the flag the next refresh would terminate an uninitialized library
and fail there, every time, until the daemon restarted.

\end{document}
